\documentclass{article}
\PassOptionsToPackage{numbers}{natbib}
\usepackage[preprint]{neurips_2025}
\usepackage[utf8]{inputenc}
\usepackage[T1]{fontenc}
\usepackage{hyperref}
\usepackage{url}
\usepackage{booktabs}
\usepackage{amsfonts}
\usepackage{nicefrac}
\usepackage{microtype}
\usepackage{xcolor}
\usepackage{listings}
\lstset{basicstyle=\ttfamily\small, breaklines=true, frame=single, columns=fullflexible}

% Every number below comes from experiments/analyze.py over committed raw logs.
% AUTO-GENERATED by experiments/analyze.py from committed raw logs. Do not edit.
\newcommand{\numntasks}{164}
\newcommand{\numpassA}{160}
\newcommand{\numpassB}{161}
\newcommand{\numpAa}{97.6}
\newcommand{\numpBb}{98.2}
\newcommand{\numciAlo}{93.9}
\newcommand{\numciAhi}{99.0}
\newcommand{\numciBlo}{94.8}
\newcommand{\numciBhi}{99.4}
\newcommand{\numdeltapp}{0.6}
\newcommand{\numbfixed}{1}
\newcommand{\numbbroke}{0}
\newcommand{\nummcnemarp}{1.0000}
\newcommand{\nummeancallsA}{1.00}
\newcommand{\nummeancallsB}{1.12}
\newcommand{\nummeansecsA}{11.6}
\newcommand{\nummeansecsB}{19.6}
\newcommand{\numtotalcallsB}{184}
\newcommand{\numfirstpassB}{146}
\newcommand{\numrepairedB}{17}
\newcommand{\numunresolvedB}{1}
\newcommand{\numpBz}{98.2}
\newcommand{\numpassBz}{161}
\newcommand{\numrepairfixedhidden}{0}
\newcommand{\numrepairbrokehidden}{0}
\newcommand{\numrepairedfalsealarm}{15}
\newcommand{\numrepairedtruebug}{2}
\newcommand{\numunresolvedhiddenpass}{1}


\title{Deliverable Contracts for Long-Horizon Agents:\\
Gating Agent Runs with Verify--Repair Loops and Tamper-Evident Audits}

\author{%
  Zhen Li\\
  \texttt{zhen@repl.it}\\
  \texttt{github.com/zhenthebuilder/charter-gate}
}

\begin{document}

\maketitle

\begin{abstract}
Long-horizon agent tasks rarely fail loudly. They fail by \emph{declaring victory}:
the agent reports completion while deliverables are missing, malformed, or silently
dropped. We propose governing agent work the way continuous integration governs human
work --- by declaring the deliverables themselves as a machine-checkable
\emph{contract}, gating every run through a bounded verify--repair loop, and recording
each attempt in a hash-chained audit log. We implement this as \textsc{charter}, a
single-file, dependency-free CLI that wraps \emph{any} agent invocable as a shell
command. To anchor the mechanism on a public benchmark, we run a paired A/B on all 164
HumanEval problems with a small model (\texttt{claude-haiku-4-5}): the baseline arm
submits the model's first answer; the gated arm adds a contract requiring the solution
to parse, define the entry point, and pass model-authored smoke tests derived
\emph{only} from the public docstring, with at most two repair rounds. Hidden tests are
never shown. Scored with the official HumanEval harness, accuracy is at ceiling and
the gate's effect on pass@1 is small and not statistically significant
(\numpAa\%\,$\rightarrow$\,\numpBb\%, +\numdeltapp{}pp, exact McNemar
$p=\nummcnemarp$); we report this null result as such. What the run does establish
is the gate's governance behaviour at low cost: it caught first-attempt contract
breaches on $\numrepairedB+\numunresolvedB$ of \numntasks{} tasks, auto-repaired
\numrepairedB{} of them, surfaced the remaining \numunresolvedB{} as an explicit
nonzero-exit failure instead of a silent claim of completion, with \numbbroke{}
regressions observed, at $\nummeancallsB\times$ the model calls of the baseline. We release the tool, the full raw logs, and the scoring pipeline; every
number in this paper is generated directly from the committed artifacts.
\end{abstract}

\section{Introduction}

Agents that work for minutes to hours produce \emph{deliverables}: files, patches,
reports, passing test suites. The dominant failure mode in practice is not a crash but
an unverified claim of completion --- the agent ``declares victory'' on deliverables
that do not exist, do not pass, or quietly omit part of the task. Current mitigations
are either \emph{post hoc} (offline evals, human review) or \emph{framework-bound}
(bespoke verifier hooks inside a particular agent stack).

Software engineering solved the same problem for humans decades ago: continuous
integration. Nobody merges on the author's word; a declarative pipeline checks the
deliverables and blocks on failure. We port this pattern to agent runs with three
components:

\begin{enumerate}
\item \textbf{Contract.} A \texttt{charter.toml} file declares deliverables (paths and
globs, minimum sizes and counts, required/forbidden regular expressions, freshness)
and arbitrary executable acceptance checks (any command; exit 0 means pass).
\item \textbf{Gate.} \texttt{charter run -- <agent command>} executes the agent,
verifies the contract, and on breach re-invokes the agent with the structured
violation report injected (via an environment variable pointing at the report), up to
a bounded number of repair rounds.
\item \textbf{Audit.} Every attempt appends a record to \texttt{.charter/audit.jsonl}
containing the verdict and the SHA-256 of every deliverable, hash-chained to the
previous record so after-the-fact tampering is detectable.
\end{enumerate}

The design is deliberately minimal: a single stdlib-only Python file, agnostic to the
agent (anything that is a shell command), installable by a stranger in seconds. Our
claim is not that contracts make a weak model strong; it is that an \emph{external,
declarative} definition of ``done'' --- checked, retried, and logged outside the
agent's own control --- converts silent victory declarations into either repaired
deliverables or explicit, auditable breaches, and yields an evidence trail that an
agent's self-report cannot provide. We evaluate both whether the gate is \emph{safe}
(does it harm benchmark accuracy?) and what its governance loop actually does,
reporting the accuracy comparison honestly as a null result at ceiling.

\paragraph{Contributions.}
(i) A problem framing: deliverable-contract governance for long-horizon agent runs,
with contract, gate, and audit as separable components.
(ii) \textsc{charter}, an open-source, dependency-free reference implementation.
(iii) A paired A/B on a public benchmark (HumanEval, official scorer, $n=\numntasks$
per arm) with a repair-loop ablation, withholding hidden tests from the model, with
all raw logs and scoring code released --- including an honest null result on
accuracy and a decomposition of what the gate's repairs actually changed.

\section{Related Work}

\textbf{Self-correction and self-refinement.} Prompting an LLM to critique and revise
its own output (e.g., Self-Refine~\cite{madaan2023selfrefine},
Reflexion~\cite{shinn2023reflexion}) improves some tasks, but the critic shares the
generator's blind spots, and unguided self-correction can degrade
performance~\cite{huang2024large}. Our gate differs in that the retry signal is an
\emph{externally executed} contract verdict (parse checks, executed tests), not the
model's opinion of itself.

\textbf{Execution-guided code generation.} Using executable tests at generation time
is well established: CodeT generates tests to select samples~\cite{chen2022codet},
LEVER verifies with execution~\cite{ni2023lever}, and self-debugging feeds execution
errors back to the model~\cite{chen2023selfdebug}. Our HumanEval arm B is an instance
of this family; the contribution is not the specific signal but its packaging as a
general, declarative, agent-agnostic contract that also covers non-code deliverables
and produces an audit chain.

\textbf{Agent evaluation harnesses.} SWE-bench~\cite{jimenez2024swebench},
GAIA~\cite{mialon2023gaia}, and HumanEval~\cite{chen2021humaneval} score agent or
model outputs offline against hidden references. Contracts are complementary: they
run \emph{inline}, use only information the operator legitimately has (the spec), and
gate the run rather than grade it afterwards.

\textbf{Policy-as-code.} Declarative gates over artifacts are standard in
infrastructure (OPA, branch protection, pre-commit). \textsc{charter} transplants
that pattern to agent deliverables and adds the bounded repair loop and hash-chained
attempt log.

\section{The \textsc{charter} system}

\subsection{Contract language}
A contract is a TOML file with two block types. \texttt{[[deliverable]]} declares a
path or glob with optional \texttt{min\_count}, \texttt{min\_bytes},
\texttt{must\_match} / \texttt{must\_not\_match} regular expressions, and
\texttt{max\_age\_minutes} (freshness). \texttt{[[check]]} declares an arbitrary
command with a timeout; exit status 0 is a pass. \texttt{charter verify} evaluates
all blocks, prints a human-readable and (optionally) JSON report, and exits nonzero
on any breach --- making it directly usable as a CI step.

\subsection{Verify--repair gate}
\texttt{charter run --max-repairs $k$ -- <agent cmd>} executes the agent in the
workspace, then verifies. On breach, it writes the violation report to a file, sets
\texttt{CHARTER\_VIOLATIONS} to its path, and re-invokes the same agent command; a
contract-aware driver includes the report (and its previous outputs) in the next
prompt. The loop is bounded by $k$, so a contract that the agent cannot satisfy
terminates with a nonzero exit rather than an infinite loop.

\subsection{Tamper-evident audit}
Each attempt appends \texttt{\{attempt, agent\_exit, ok, deliverable\_sha256s, prev\}}
to \texttt{.charter/audit.jsonl}, where \texttt{prev} is the SHA-256 of the previous
record; \texttt{charter log} re-verifies the chain. The operator can therefore prove
which bytes the agent delivered on which attempt, and detect post-hoc modification of
the log itself.

\section{Experimental setup}

\paragraph{Benchmark and model.} We use all \numntasks{} problems of
HumanEval~\cite{chen2021humaneval} and score with the official open-source harness
(vendored unmodified; the one documented line enabling execution is uncommented).
The subject model is \texttt{claude-haiku-4-5} invoked through the \texttt{claude}
CLI; we chose a small model because frontier models are near ceiling on HumanEval,
which would make any gate effect unmeasurable. HumanEval is not itself long-horizon;
we use it because its unit of evaluation \emph{is} a deliverable with an acceptance
test, the property the gate targets, and because it admits official, local,
reproducible scoring within our compute budget (Section~\ref{sec:limitations}
discusses external validity).

\paragraph{Conditions (paired, per task).}
\textbf{A (baseline):} one model call; the produced \texttt{solution.py} is submitted
as-is. \textbf{B (gated):} the same model and task prompt, but the model must produce
both \texttt{solution.py} and \texttt{smoke\_test.py}, a self-contained test derived
\emph{only} from the public docstring. A per-task contract requires: the solution
file exists with the entry-point definition; it parses (\texttt{ast.parse}); the smoke
test exists, imports the solution, and passes when executed. The gate runs with
\texttt{--max-repairs 2}, so at most three model calls per task. \emph{The hidden
HumanEval tests are never shown to the model or the contract in either arm.} The
extra information available to arm B is exactly what an operator has: the spec and
the ability to execute what the agent wrote.

\paragraph{Scoring and statistics.} Official harness pass@1 per arm; Wilson 95\%
confidence intervals; and, because the design is paired, an exact McNemar test
(binomial test on discordant pairs). We additionally score \textbf{B$_0$}, arm B's
first-attempt solutions re-extracted from the committed raw model outputs, which
ablates the repair loop and separates the effect of the arm-B prompt (write tests
alongside the solution) from the effect of verify--repair. All raw model outputs
(\texttt{call\_n.txt}),
per-task metadata, charter audit logs, official per-task verdicts, and the analysis
script are committed; \texttt{experiments/analyze.py} regenerates every number in
this section, including this paper's macro file.

\section{Results}

\begin{table}[h]
\centering
\caption{Paired A/B on HumanEval ($n=\numntasks$ per arm), official scorer,
\texttt{claude-haiku-4-5}. Arm B adds the \textsc{charter} gate (docstring-derived
smoke tests, $\le 2$ repairs); hidden tests withheld in both arms. B$_0$ scores arm
B's \emph{first-attempt} solutions (extracted from the committed raw outputs), i.e.
arm B with the repair loop ablated.}
\label{tab:main}
\begin{tabular}{lcccc}
\toprule
Arm & pass@1 & 95\% CI & passed/$n$ & mean model calls \\
\midrule
A: single shot          & \numpAa\% & [\numciAlo, \numciAhi] & \numpassA/\numntasks & \nummeancallsA \\
B$_0$: B without repairs & \numpBz\% & --- & \numpassBz/\numntasks & \nummeancallsA \\
B: \textsc{charter} gate & \numpBb\% & [\numciBlo, \numciBhi] & \numpassB/\numntasks & \nummeancallsB \\
\bottomrule
\end{tabular}
\end{table}

\paragraph{Accuracy: a null result at ceiling.} The baseline already passes
\numpassA{} of \numntasks{} hidden test suites (\numpAa\%), leaving headroom of only
four tasks. The gated arm passes \numpassB{} (+\numdeltapp{}pp). The paired
decomposition shows \numbfixed{} task fixed and \numbbroke{} broken (exact McNemar
$p=\nummcnemarp$): \emph{the accuracy difference is not statistically significant},
and we report it as a null result. The B$_0$ ablation localises the single fixed task
(HumanEval/92) to the arm-B \emph{prompt} (writing acceptance tests alongside the
solution), not the repair loop: repairs changed \numrepairfixedhidden{} hidden-test
outcomes (\numrepairbrokehidden{} regressions). A gate cannot add correctness that a
near-ceiling model does not lack; measuring an accuracy effect would require a harder
benchmark or weaker model, which we flag as future work rather than re-rolling the
dice here.

\paragraph{Governance behaviour: what the gate actually did.} Charter's audit chains
decompose arm B: \numfirstpassB{} tasks satisfied the contract on the first attempt;
\numrepairedB{} breached it and were auto-repaired within two rounds;
\numunresolvedB{} ended with an unresolved breach, which the gate surfaced as a
nonzero exit --- the operator-visible alternative to a silent victory declaration.
Among the \numrepairedB{} repaired tasks, \numrepairedfalsealarm{} were
\emph{false alarms} (the first-attempt solution already passed the hidden tests; the
model's own smoke test was wrong, and the repair round fixed the test), while
\numrepairedtruebug{} were genuine solution bugs that the bounded repair loop did not
fix against hidden tests. The unresolved task's solution in fact passes the hidden
tests (\numunresolvedhiddenpass/\numunresolvedB): the gate's verdicts are exactly as
good as the contract, no better. In this run we observed \numbbroke{} regressions
relative to baseline; with a single run and one discordant pair this is evidence of
absence of harm only at the run level, not a general safety guarantee.

\paragraph{Cost.} The gate cost $\nummeancallsB\times$ model calls
(\numtotalcallsB{} total vs.\ 164) and raised mean per-task model latency from
\nummeansecsA{}s to \nummeansecsB{}s --- the price of executing smoke tests and up to
two repairs on the \numrepairedB$+$\numunresolvedB{} breaching tasks.

\paragraph{Honest reading.} On this benchmark, with this model, the verify--repair
gate bought no significant accuracy and mostly adjudicated the model's own buggy
self-tests. What it provably delivered --- visible in the committed audit chains ---
is the governance contract: every task ended either contract-satisfied or with an
explicit, machine-readable breach, with zero regressions and a 12\% call overhead.
We consider publishing the null accuracy result alongside the mechanism more useful
to practitioners than a benchmark cherry-picked to flatter the gate.

\section{Dogfooding: governing this project's own deliverables}
The repository's root \texttt{charter.toml} declares this project's five contracted
deliverables (deployable CLI, this paper compiled from auto-generated numbers, the
benchmark's raw logs and official scores, a marketing plan, and a static landing
page) with size, pattern, and count constraints plus executable checks --- including
one that re-runs the analysis script. \texttt{charter verify} passes on the released
revision; the audit chain for the benchmark runs is included. This is anecdotal
rather than experimental evidence, but it demonstrates the intended long-horizon use:
the contract was written before the deliverables existed and gated the run that
produced them.

\section{Limitations}
\label{sec:limitations}

\textbf{Ceiling effect / null accuracy result.} \texttt{claude-haiku-4-5} passes
\numpAa\% of HumanEval single-shot, leaving four tasks of headroom; no gate could
show a meaningful accuracy effect in this regime, and ours did not. The accuracy
claim of this paper is therefore only ``\numbbroke{} regressions were observed in
this run'' --- which does not establish harmlessness in general, given one run and
one discordant pair; any positive accuracy claim needs a harder benchmark (SWE-bench Lite
with a full scaffold is the natural next step) or a weaker model.
\textbf{Self-verification quality.} The gate is only as good as the contract. Here
\numrepairedfalsealarm{} of \numrepairedB{} repairs fixed the model's own buggy
smoke test rather than the solution, and the one unresolved breach was itself a
false alarm; docstring examples also partially overlap hidden tests. Operators who
write their own acceptance checks (the intended deployment) sidestep the
self-authored-test failure mode but inherit contract-authoring cost.
\textbf{External validity.} HumanEval is short-horizon and single-file; it anchors
the gate's mechanics on a public, officially-scored benchmark, not the full
long-horizon claim. \textbf{Single model, single benchmark, $n=1$ run per arm; no
variance estimate across reruns.} \textbf{Scope.} Contracts check artifacts, not
intent; a malicious agent could satisfy the letter of a weak contract. The audit
chain is tamper-\emph{evident}, not tamper-\emph{proof}.

\section{Conclusion}
Treating ``done'' as a contract --- declared outside the agent, verified by
execution, retried under a bound, and logged in a tamper-evident chain --- is a
cheap, agent-agnostic way to govern the deliverables of long-horizon agent runs. On
a public benchmark at ceiling, the gate changed accuracy by a statistically
insignificant amount (reported here as a null result) while doing what governance
tooling must do: it converted every run into either a contract-satisfied deliverable
or an explicit, auditable breach, with zero regressions observed in this run, at
12\% extra model calls.
The tool, the full raw logs, and the analysis that generated every number above are
public: \url{https://github.com/zhenthebuilder/charter-gate}.

\begin{thebibliography}{9}

\bibitem{chen2021humaneval}
M. Chen et al.
\newblock Evaluating large language models trained on code.
\newblock \emph{arXiv:2107.03374}, 2021.

\bibitem{madaan2023selfrefine}
A. Madaan et al.
\newblock Self-Refine: Iterative refinement with self-feedback.
\newblock In \emph{NeurIPS}, 2023.

\bibitem{shinn2023reflexion}
N. Shinn et al.
\newblock Reflexion: Language agents with verbal reinforcement learning.
\newblock In \emph{NeurIPS}, 2023.

\bibitem{huang2024large}
J. Huang et al.
\newblock Large language models cannot self-correct reasoning yet.
\newblock In \emph{ICLR}, 2024.

\bibitem{chen2022codet}
B. Chen et al.
\newblock CodeT: Code generation with generated tests.
\newblock In \emph{ICLR}, 2023.

\bibitem{ni2023lever}
A. Ni et al.
\newblock LEVER: Learning to verify language-to-code generation with execution.
\newblock In \emph{ICML}, 2023.

\bibitem{chen2023selfdebug}
X. Chen, M. Lin, N. Sch\"arli, and D. Zhou.
\newblock Teaching large language models to self-debug.
\newblock In \emph{ICLR}, 2024.

\bibitem{jimenez2024swebench}
C. Jimenez et al.
\newblock SWE-bench: Can language models resolve real-world GitHub issues?
\newblock In \emph{ICLR}, 2024.

\bibitem{mialon2023gaia}
G. Mialon et al.
\newblock GAIA: A benchmark for general AI assistants.
\newblock In \emph{ICLR}, 2024.

\end{thebibliography}

\end{document}
